\section*{الفصل \ref{ch:b2:continuous-reactors} --- المفاعلات المستمرة والعمليات الصناعية}

\begin{solution}{exo:b2:continuous-reactors:1}
$\tau = V/Q_v = 2000/0.50 = \qty{4.0e3}{s}$، أي \qty{1.1}{h}.
\end{solution}

\begin{solution}{exo:b2:continuous-reactors:2}
$k\tau = 1.2 \times 10^{-3} \times 1800 = 2.16$. الخزان: $X = 2.16/3.16 = 0.68$؛
الأنبوب: $X = 1 - \mathrm e^{-2.16} = 0.88$.
\end{solution}

\begin{solution}{exo:b2:continuous-reactors:3}
الأنبوب: $\tau = \ln 100/0.020 = \qty{230}{s}$. الخزان: $X/(1 - X) = 99 = k\tau$،
$\tau = \qty{4950}{s}$، أي أطول بأكثر من عشرين مرة: فعند \omterm{def:b2:continuous-reactors:conversion}{تحويل} مرتفع يدفع
الخزان المحرَّك ثمنًا باهظًا لعمله عند تركيز المخرج.
\end{solution}

\begin{solution}{exo:b2:continuous-reactors:4}
لكل لتر: يتحرر $120 \times 0.50 = \qty{60}{kJ}$، فيسخّن
\qty{4.2}{kJ/K}: $\Delta T_{\text{كظوم}} = \qty{14}{K}$. وتعطل التبريد
يدفئ الخليط بوضوح لكن دون خطر، ما لم يكن المذيب
قريبًا من نقطة غليانه أو يبدأ تفاعل ثانٍ في هذا المجال.
\end{solution}

\begin{solution}{exo:b2:continuous-reactors:5}
الخزان: $Da = X/(1 - X)^2 = 0.8/0.04 = 20$، $\tau = 20/0.010 = \qty{2000}{s}$.
الأنبوب: $Da = X/(1 - X) = 4$، $\tau = \qty{400}{s}$. النسبة 5، مقابل $4/\ln 5 =
2.5$ للرتبة 1: فكلما ارتفعت الرتبة ازداد ما يعانيه الخزان المحرَّك
من تركيزه المنخفض.
\end{solution}

\begin{solution}{exo:b2:continuous-reactors:6}
$k\tau = 5.0$. ثلاثة خزانات: $X = 1 - (1 + 5/3)^{-3} = 0.947$. خزان واحد: $5/6 =
0.833$. الأنبوب: $1 - \mathrm e^{-5} = 0.993$.
\end{solution}

\begin{solution}{exo:b2:continuous-reactors:7}
$X = 0.70$ يعني $k\tau = 0.70/0.30 = 2.33$. ومضاعفة التدفق تنصّف $\tau$:
$k\tau = 1.17$، $X = 0.54$. ولاسترجاع 70\,\% يجب مضاعفة \omterm{def:b2:continuous-reactors:residence-time}{زمن المكوث}، ومن ثم
الحجم.
\end{solution}

\begin{solution}{exo:b2:continuous-reactors:8}
المحررة: $80 \times 2.0 \times 1.0 \times 0.90 = \qty{144}{kW}$. وما يحمله
التيار الخارج، المسخَّن من \qty{330}{K} إلى \qty{350}{K}: $4.0 \times 1.0 \times 20 =
\qty{80}{kW}$. وعلى الغلاف أن يصرف الفرق، \qty{64}{kW}.
\end{solution}

\begin{solution}{exo:b2:continuous-reactors:9}
الأبعاد الخطية $\times 10$: الحجم $\times 1000$، ومساحة الجدار $\times 100$، والمساحة
لكل وحدة حجم $\div 10$. فالحرارة المحررة تكبر مع الحجم، والحرارة
المصروفة عبر الجدار مع مساحته: فالمفاعل الكبير يصرف، لكل لتر،
حرارة أقل بعشر مرات عند فرق درجة الحرارة نفسه.
\end{solution}

\begin{solution}{exo:b2:continuous-reactors:10}
الأنبوب: $\tau = \ln(0.05/0.10)/(0.05 - 0.10) = \qty{13.9}{min}$، $c_B/c_{A,0} =
\frac{0.10}{-0.05}(\mathrm e^{-1.386} - \mathrm e^{-0.693}) = 0.50$. الخزان: $\tau =
1/\sqrt{0.005} = \qty{14.1}{min}$، $c_B/c_{A,0} = 1.414/(2.414 \times 1.707) =
0.34$. ففي الخزان يختلط A الطازج وB المكتمل في الحجم نفسه: فيبقى بعض
B دائمًا مدة كافية ليتحول إلى C بينما بعض A بالكاد
وصل؛ أما في الأنبوب فلكل شريحة التاريخ نفسه.
\end{solution}

\begin{solution}{exo:b2:continuous-reactors:11}
رفع درجة حرارة سائل التبريد (والتغذية) يزيح مستقيم الصرف نحو
اليمين. فترتفع نقطة التقاطع الباردة ببطء حتى يصير المستقيم مماسًا
للانحناء السفلي من حرف S؛ وبعد ذلك تختفي الحالة الباردة ويقفز الخزان
إلى الفرع الساخن (الاشتعال). وعند خفض درجة الحرارة من جديد، يبقى الخزان
على الفرع الساخن حتى يصير المستقيم مماسًا للانحناء العلوي، ولا يعود
إلا عندئذ (الانطفاء)، عند درجة حرارة أدنى من درجة الاشتعال: إنها
حلقة تخلّف.
\end{solution}

\begin{solution}{exo:b2:continuous-reactors:12}
\omterm{def:b2:continuous-reactors:reactors}{مفاعل الدفعات}: يبقى 75\,\% من \qty{2.0}{mol/L}، أي $1.5 \times 100 = \qty{150}{kJ/L}$،
و$\Delta T = 150/4.0 = \qty{37.5}{K}$ فوق درجة حرارة التشغيل، بما يكفي
لبلوغ نقطة غليان أو تفكك. والتشغيل نصف المستمر: لا يزيد غير المتفاعل
على \qty{0.10}{mol/L}، $\Delta T = 10/4.0 = \qty{2.5}{K}$. فالإضافة البطيئة
تحدّ الطاقة المختزنة في المفاعل.
\end{solution}

\begin{solution}{pb:b2:continuous-reactors:1}
\textbf{1.} $r = k[\text{إستر}][\ce{OH-}]$؛ فالتغذيتان المتساويتان والستوكيومترية 1 : 1
تحفظان تساوي التركيزين، $r = kc^2$.
\textbf{2.} $1/c - 1/c_0 = kt$ مع $c = 0.10c_0$: $kc_0t = 9$، $t = 9/(0.11
\times 0.10) = \qty{818}{s}$، أي نحو \qty{14}{min}.
\textbf{3.} $t_{1/2} = 1/(kc_0) = \qty{91}{s}$؛ ففي الرتبة 2 تنخفض السرعة مع
مربع التركيز، فيتعلق زمن تنصيفه بنقطة
الانطلاق.
\textbf{4.} $Da = kc_0\tau = 0.011\,\tau$ ($\tau$ بالثواني).
\textbf{5.} $Q_vc_0 - Q_vc - kc^2V = 0$، أي $c_0 - c = k\tau c^2$.
\textbf{6.} مع $c = c_0(1 - X)$: $c_0X = k\tau c_0^2(1 - X)^2$.
\textbf{7.} $Da = 0.90/0.10^2 = 90$.
\textbf{8.} $\tau = 90/0.011 = \qty{8.2e3}{s}$، أي \qty{2.3}{h}.
\textbf{9.} $V = 0.50 \times 8182 = \qty{4.1e3}{L}$، أي \qty{4.1}{m^3}.
\textbf{10.} تركيز المخرج، $0.010\,\unit{mol/L}$: فالسرعة في كل مكان
أدنى بمئة مرة منها عند المدخل.
\textbf{11.} $Da = X/(1 - X) = 9$، $\tau = \qty{818}{s}$، $V = \qty{409}{L}$.
\textbf{12.} حجم الأنبوب أصغر بعشر مرات.
\textbf{13.} $c_{j-1} - c_j = k\tau_1c_j^2$، حيث $\tau_1$ زمن مكوث
خزان واحد.
\textbf{14.} $\tau_1 = 850/(3 \times 0.50) = \qty{567}{s}$، $k\tau_1 =
\qty{62.4}{L/mol}$. وبحل $62.4c_j^2 + c_j - c_{j-1} = 0$: $c_1 = 0.0328$،
$c_2 = 0.0163$، $c_3 = \qty{0.0100}{mol/L}$: \omterm{def:b2:continuous-reactors:conversion}{تحويل} قدره 90\,\%.
\textbf{15.} الأنبوب، أو السلسلة (أصغر من الخزان الواحد بعشر مرات وبخمس
مرات). وقد يفوز الخزان الواحد مع ذلك بدرجة حرارته المتجانسة وسهولة
تنظيفه وضبطه، ولأن الحجم هنا رخيص.
\textbf{16.} $\Delta T_{\text{كظوم}} = 55\,000 \times 100/(997 \times 4180) =
\qty{1.3}{K}$ (\qty{1.2}{K} عند 90\,\%).
\textbf{17.} $55 \times 0.50 \times 0.10 \times 0.90 = \qty{2.5}{kW}$.
\textbf{18.} لا: يحمل التيار الخارج الحرارة بارتفاع نحو
كلفن واحد؛ فلا جموح ممكن.
\textbf{19.} تركيز أعلى بعشرين مرة: يصير الارتفاع الكظوم
\qty{26}{K} والحرارة \qty{50}{kW}، وهو ما يستحق ملف تبريد؛ ولأن المقدار $kc_0$
أكبر بعشرين مرة، يصير الخزان اللازم \omterm{def:b2:continuous-reactors:conversion}{لتحويل} 90\,\% أصغر بعشرين مرة
(\qty{0.20}{m^3}).
\textbf{20.} $0.50 \times 0.10 \times 0.90 = \qty{0.045}{mol/s}$، $\times 86\,400
\times 82.03 = \qty{319}{kg}$ من إيثانوات الصوديوم في اليوم.
\textbf{21.} يكبر ثابت السرعة مع درجة الحرارة (أرهينيوس)، فيصغر
الحجم؛ والثمن تسخين التغذية، وفي تفاعلات أخرى \omterm{def:b2:continuous-reactors:selectivity}{انتقائية}
أدنى وبخار أكثر.
\textbf{22.} يجب أن يسع الخزان المحرَّك الواحد $\boldsymbol{V \approx
\qty{4.1}{m^3}}$.
\end{solution}
